\chapter{Introducción}

\section{Motivación}

El origen de esta tesis no es un algoritmo, sino una pregunta que arrastro desde hace 
tiempo: ¿se podría mostrar en una pantalla lo que alguien imagina? Al oír la palabra 
``imaginación'' uno piensa enseguida en entretenimiento, pero el asunto va mucho más 
lejos. El desarrollo de esta tecnología podría ayudarnos a entender qué sucede en la 
mente de personas con esquizofrenia, o en general cualquier transtorno correspondiente 
a la alteracion en la psique y corteza cerebral, lugar donde la frontera entre lo que 
se ve y lo que se imagina se vuelve borrosa. Además de abrir un canal de comunicación 
para pacientes con síndrome de enclaustramiento, que piensan con normalidad aunque no 
logran moverse. Y podría dar voz a personas con paraplejia, hoy limitadas a la mirada 
o a unos pocos gestos. 

Hoy nos entendemos con las máquinas escribiendo, hablando o haciendo clic; o siendo 
mas tecnicos por dispositivos de interaccion humano-computadora, en ambos casos pasa 
por las manos y por un control voluntario. Una interfaz que leyera directamente el 
contenido mental se saltaría ese cuello de botella. El problema es que descifrar 
imágenes desde el cerebro cuesta muchísimo en términos de cómputo, debido a los sitemas
de señaes como las EEG no tienen a ser muy precisos, los sistemas fNIRS de espectrocopia
de infrarrojo cercano, si bien son mas precisos, no es tan potente como una fMRI;
por otro lado la alta dimensionalidad del cerebro deja pequeño a cualquier conjunto de 
datos, de ahí que este trabajo todavía no se meta con la imaginería. 

Lo primero es reconstruir lo que el sujeto está viendo con su fMRI; el salto a lo que
imagina vendrá después. El objetivo de fondo, mirando a varios años adelante, es 
proyectar directamente la imaginacion a entornos digitales.

\section{Planteamiento técnico}

Decodificar imágenes desde fMRI es una especie de problema inverso mal 
condicionado. La señal depediente del nivel de oxigenacion sanguiena(BOLD),del conjunto 
de datos BOLD5000, aporta del orden de $V \approx 50{,}000$ vóxeles frente a unos 
escasos $N \approx 5{,}000$ estímulos; y, para colmo, la respuesta hemodinámica no es 
igual de una persona a otra \cite{logothetis2008what}. Shen et al.~\cite{shen2019deep} 
y Beliy et al.~\cite{beliy2019voxels} mostraron que las primeras capas de VGG-19 
explican parte de la actividad de la corteza visual, trabajo llevado a cabo con arquitectura VQGAN. 
Pero dejaron algo igual de importante, el querer reconstruir la imagen original sin 
pasar por un nivel semántico termina en imágenes borrosas, sin identidad de objeto.

Con la llegada de CLIP \cite{radford2021learning} y los modelos de difusión latente 
\cite{rombach2022high} el planteamiento cambió de raíz. Ya no hace falta regresar a los
píxeles, basta con apuntar al espacio de embeddings y dejar que el generador arme la 
imagen. Partiendo de esa idea monté un primer prototipo, el adapter lineal Ridge que 
lleva de fMRI a CLIP ViT-L/14 y pasa el resultado a Stable Diffusion 2.1 unCLIP. 
Opte por lo lineal porque mostraba hasta dónde llegaba un mapeo proyectivo y ayudaba a 
ubicar en qué punto se rompe la decodificación.

Para tener un baseline matemáticamente claro antes de complicar nada, la primera fase 
implementa ese proyector Ridge, que estabiliza el problema inverso con regularización 
de Tikhonov \cite{tikhonov1977solutions}:
\begin{equation}
\hat{\beta} = \arg\min_{\beta} \| Y - X\beta \|_2^2 + \lambda \|\beta\|_2^2
\end{equation}
donde $X$ es la matriz de diseño con los vóxeles e $Y$ la matriz de embeddings CLIP 
objetivo. Insisto en su papel: el Ridge Lineal está aquí para diagnosticar, no para 
competir. Deja ver el efecto sin que la no linealidad o el ruido lo tapen.

\section{Problema}

El diseño actual abarca una dificultad que aparece una y otra vez; CLIP y fMRI no son 
isomorfos. El adapter puede generar vectores en $\mathbb{R}^{768}$, con tan altisima 
dimensionalidad que no hay garantía que se encuentren dentro del colector de 
representaciones de  CLIP. Entonces \textbf{¿Cómo bajar la alietoriedad en la 
reinterpretacion de los embeddings de fMRI de BOLD5000 por parte de CLIP mientras se mantiene la 
coherencia global de la imagen reconstruida?}

\section{Objetivo General}

Armar y diagnosticar un sistema que reconstruye imagen a partir de fMRI sobre BOLD5000, 
usando dos adapters. El primero, un proyector lineal Ridge, sirve de línea base 
bien aislable. El segundo, un adapter Ridge Estocástico inspirado en MindEye y 
MindEye2 que arregla el comportamiento determinista del lineal y la pérdida de norma 
en el espacio CLIP. Los dos alimentan el pipeline image-to-image de Stable Diffusion 2.1 
unCLIP.

\section{Objetivos Específicos}

\begin{itemize}
\item[•] Preprocesar BOLD5000 y armar la máscara cortical visual (V1--V4, LOC, PPA), 
cambiando la selección heurística por una restricción biológica.
\item[•] Ajustar el adapter Ridge y barrer el hiperparámetro $\lambda$ con validación 
cruzada en hojas por sujeto.
\item[•] Programar el Ridge Estocástico, que cuesta casi lo mismo que el lineal salvo 
por el término de ruido que se inyecta sobre la dirección predicha.
\item[•] Enganchar ambos adapters al pipeline SD 2.1 unCLIP y calibrar el factor de guía 
CFG \cite{ho2022classifier}.
\item[•] Medir con SSIM, PSNR, PixCorr, LPIPS, CLIP-cos y 2AFC, y leer los modos de 
fallo de forma cualitativa.
\end{itemize}

\section{Hipótesis}

\begin{itemize}
\item[H1.] El Ridge capta parte de la dirección semántica, pero ni respeta la norma ni 
la geometría angular de CLIP. Esa carencia explica las fallas tipográficas que vi en mis 
ensayos previos, las mismas que el Modality Gap anticipa en el plano teórico 
\cite{liang2022mind}.
\item[H2.] Al inyectar variabilidad controlada sobre la salida del proyector, el Ridge 
Estocástico corre las predicciones hacia zonas del manifold CLIP donde SD 2.1 unCLIP las 
procesa mejor. La consecuencia esperada es más coherencia semántica en CLIP-cos, y todo 
ello sin pagar un entrenamiento contrastivo completo.
\end{itemize}

\section{Justificación}

Este estudio propone un puente interdisciplinario entre la neurociencia cognitiva y la in
teligencia artificial. Busca llenar un vacío en el conocimiento al explorar cómo los modelos
computacionales de semántica pueden alinearse con los mapas semánticos unificados del cerebro
humano y la codificación de información relacional. Al hacerlo, se sientan las bases teóricas
para la traducción computacional de la semántica visual, avanzando en la comprensión de
cómo la información se representa tanto en sistemas biológicos como artificiales. 

Las aplicaciones potenciales de esta tecnología son transformadoras y de alto impacto
social. En el ámbito clínico, podría ofrecer vías de comunicación avanzadas para pacientes
con discapacidades severas, permitiéndoles externalizar pensamientos visuales complejos.
En la neurociencia, serviría como una herramienta computacional para validar modelos de
codificación cerebral, afinar la decodificación intermodal abre la puerta, más adelante, a BCIs 
no invasivas para pacientes con síndrome de enclaustramiento, a estudios de patologías donde 
se mezclan lo percibido y lo imaginado, y a canales de salida para personas con discapacidad 
motora grave permitiendo a los investigadores ”ver” lo que el modelo predice que el cerebro está 
representando.

El enfoque de esta investigación es novedoso, ya que explora la síntesis generativa condicionada 
por la actividad cerebral, alejándose de la simple clasificación de categorías. Propone la integración 
de arquitecturas de vanguardia (LDMs y CLIP) que han demostrado resultados de alta fidelidad en los 
estudios más recientes, aplicando estos métodos a la decodificación de señales complejas. 
El diseño conceptual resultante servirá de guía metodológica para futuras implementaciones 
experimentales en el laboratorio.

\section{Estructura de la Tesis}

\begin{itemize}
\item Cap. 2 --- Marco Teórico: Este capítulo arma las bases teóricas que se necesitan para 
entender el problema de traducir entre señales cerebrales y espacios latentes visuales.
Está dividido en cuatro bloques que se complementan: la biología de la
señal neuronal (Neurociencia Computacional), cómo se formaliza el mapeo
en alta dimensionalidad (Matemáticas y Computación), las representaciones
semánticas (Inteligencia Artificial Multimodal) y los motores que arman
las imágenes (Modelos Generativos Profundos)

\item Cap. 3 --- Estado del Arte: Lo que arrancó como un intento de reconstruir
píxeles directo hoy se mira como un problema de traducción semántica entre
dos mundos: el del córtex y el de los espacios latentes multimodales. Este
capítulo cuenta esa evolución desde las primeras redes convolucionales hasta
lo que se hace ahora con difusión y aprendizaje contrastivo.

\item Cap. 4 --- Herramientas: Este capítulo cuenta qué librerías de procesamiento de fMRI, 
modelos preentrenados, lenguajes y recursos de cómputo se usaron en su entorno además de
arquitecturas propias desarrolladas para el proyecto.

\item Cap. 5 --- Análisis : Este capítulo cuenta cómo está armado BOLD5000, el preprocesamiento, 
los pares (estímulo, beta), las particiones, el modelo entidad-relación, el prototipeo, una 
mirada a los algoritmos existentes y a los propios, y los requerimientos del sistema.

\item Cap. 6 --- Desarrollo de la Investigación: Este capítulo formaliza los dos adapters fMRI→CLIP que se usan en 
esta tesis, cuenta cómo se integran con Stable Diffusion 2.1 unCLIP, detalla el procedimiento 
de inferencia y deja documentado cómo se arma cada métrica de evaluación. 

\item Cap. 7 --- Análisis de Resultados: Este capítulo cuenta cómo va la evaluación. 
La estrategia experimental tiene dos fases lógicas: (i) el adapter Ridge Lineal nativo a 
CLIP ViT-L/14 con los bugs arreglados, que sirve de diagnóstico cuantitativo y cualitativo 
del fallo del mapeo lineal puro sobre BOLD5000; y (ii) el adapter Ridge Estocástico, que le 
suma ruido Gaussiano calibrado al predictor lineal y renormaliza el embedding antes de mandarlo 
al pipeline. La corrida cuantitativa final del Ridge Estocástico se hace en la GPU NVIDIA RTX 
4070 Ti, siguiendo la regla del flujo de trabajo. Aquí se reportan métricas, comparativos y 
ejemplos visuales de la fase (i), y se dejan marcadas las tablas que se completarán con los 
resultados de la fase (ii).

\item Cap. 8 --- Conclusiones: En este capitulo se recapitulan los hallazgos del 
estudio, tanto tecnicos y cientificos del trabajo, comparando las hipotesis del Cap. 1, 
con los resultados obtenidos en el Cap. 7. Sintetizando los aportes hechos.

\item Cap. 9 --- Trabajo Futuro: se trazan 3 futuras lineas de investigacion, para 
profundizar en el trabajo, como vendria a ser el pulir el Pipeline actual, usar otras fuentes
de material fMRIs o fNIRs como las de OpenNeuro y NSD para armar ensemblers a mediano plazo
y por ultimo, precisar en una decodificación instantanea con BCI multimodales que combinen
EEG y fNIRS, para la reconstruccion de imagen en tiempo real y dejar formalizado el problema 
como obtimizacion P-NP sobre flujos en vivo.
\end{itemize}

\afterpage{\blankpage}

